% ECONOMICS_PROBLEM: {"id": "C72-656", "title": "Threshold Nash equilibrium in extensive form: resolved classification", "jel_code": "C72", "source_id": "OP-0142"}
% !TeX program = xelatex
\documentclass[11pt,a4paper]{article}
\usepackage[margin=23mm]{geometry}
\usepackage{fontspec}
\setmainfont{Latin Modern Roman}
\usepackage{amsmath,amssymb,mathtools}
\usepackage{xcolor,enumitem,booktabs,longtable,array,xurl,hyperref}
\definecolor{muted}{HTML}{53636C}
\hypersetup{colorlinks=true,urlcolor=blue,linkcolor=blue}
\setlength{\parindent}{0pt}
\setlength{\parskip}{5pt}
\newcommand{\R}{\mathbb R}
\newcommand{\E}{\mathbb E}
\newcommand{\N}{\mathbb N}
\newcommand{\ind}{\mathbf 1}
\newcommand{\Var}{\operatorname{Var}}
\newcommand{\Cov}{\operatorname{Cov}}
\newcommand{\supp}{\operatorname{supp}}
\DeclareMathOperator*{\argmin}{arg\,min}
\DeclareMathOperator*{\argmax}{arg\,max}
\newcommand{\entrymeta}[3]{{\sffamily\small\color{muted}\textbf{\texttt{#1}}\quad|\quad #2\par #3\par}}
\newcommand{\Problemref}[1]{\texttt{#1}}
\newcommand{\JELref}[2]{\texttt{#2} (JEL \texttt{#1})}
\begin{document}
\section*{Threshold Nash equilibrium in extensive form: resolved classification}
\textbf{Problem ID:} \texttt{C72-656}\quad \textbf{JEL:} \texttt{C72}\par
% BEGIN ECONOMICS STATEMENT
\label{problem:OP-0142}
% GAME_THEORY_PROBLEM: {"local_id": "GT-012", "original_number": 12, "kind": "H", "entry_kind": "statement_only", "published": false, "review_required": true, "category": "game_theory"}
\label{game:GT-012}
{\sffamily\small\color{muted}
\textbf{GT-012} | Extensive-form equilibrium selection\\
\textbf{H} | General classification resolved; vector version is an elementary formulation variant.
\par}
\subsection*{Definitions and mathematical statement}
Let $G$ be an explicitly listed finite perfect-recall extensive-form game with rational chance probabilities and terminal utilities, and with the number of players included in the input. Let $b$ range over behavioral strategy profiles and $U_i(b)$ denote expected utility. Given a subset $J$ of the players and rational thresholds $(t_i)_{i\in J}$, the decision language is
\[
 \left\{(G,J,t):\exists b\ 
 \begin{array}{l}
 U_i(b)\geq U_i(b_i',b_{-i})\quad\text{for every }i,b_i',\\
 U_j(b)\geq t_j\quad\text{for every }j\in J
 \end{array}\right\}.
\]
The scalar version in [S1] uses an explicitly specified rational linear objective $\omega(U)=\sum_iw_iU_i$ and asks whether an equilibrium satisfies $\omega(U(b))\geq\lambda$. It is $\exists\mathbb R$-complete. The displayed vector version lies in the same class: membership adds finitely many polynomial inequalities; scalar-objective hardness transfers by adding a passive player whose terminal payoff is $\sum_iw_iu_i$ and giving only that player a threshold.
Here $\exists\mathbb R$ consists of decision problems polynomial-time reducible to feasibility of finite systems of polynomial equalities and inequalities over the real numbers.
\subsection*{Short English statement}
Deciding whether some equilibrium achieves specified payoff targets has existential-real complexity.
\subsection*{Variants and scope boundaries}
This is a general multiplayer result, with chance allowed. It does not establish identical hardness for every fixed-player or chance-free subclass. Strict thresholds, optimization to prescribed additive accuracy, and rational-only equilibrium witnesses are different formulations.
\subsection*{Source relationship and status}
Theorem 2 and Section 6.1 of the 2026 version resolve the scalar problem. The passive-player reduction above is an explicitly identified editorial deduction for the original catalog's vector statement; it is not claimed as verbatim source text.
\subsection*{References}
\begingroup\footnotesize\raggedright
{\footnotesize\raggedright
\textbf{[S1]} Vincent Cheval, Florian Horn, Soumyajit Paul, and Mahsa Shirmohammadi (2026). \emph{On the Complexity of the Optimal Correlated Equilibria in Extensive-Form Games}. arXiv:2507.11509v2. Locator: Theorem 2; Section 6.1. \url{https://arxiv.org/pdf/2507.11509v2}\par}
\par\endgroup

\subsection*{Common conventions: Source mathematical conventions}

\textbf{Finite games.} For a finite set $A$, $\Delta(A)$ is its probability
simplex. Players choose independent mixed strategies in Nash equilibrium;
correlation is allowed only when explicitly part of the solution concept.
In a game with payoff functions $u_i$ and strategy profile $x$, an additive
$\varepsilon$ Nash equilibrium satisfies
\[
 u_i(x)\geq u_i(x_i',x_{-i})-\varepsilon
 \quad\text{for every player $i$ and every admissible deviation $x_i'$}.
\]
For numerical approximation, payoffs are normalized as locally specified.
An $\varepsilon$ payoff residual is distinct from distance at most
$\varepsilon$ to the set of exact equilibria. Point convergence, convergence
of empirical play, and convergence in distance to an equilibrium set are
also distinct.

\textbf{Computational representations.} Unless stated otherwise, a complexity
question takes rational data with binary encoding; polynomial time is measured
in total bit length. A fixed-accuracy polynomial algorithm is not necessarily
polynomial in $1/\varepsilon$, and neither is necessarily polynomial in
$\log(1/\varepsilon)$. Succinct games and value, demand, or sampling oracles
must declare their interfaces. Exact real solutions require an explicit
representation; an unspecified list of arbitrary real numbers is not a
finite computational output. A claimed algorithmic barrier is meaningful
only for its specified model and reductions.

\textbf{Dynamic games.} Histories, information, observation, and allowable
randomization are part of the model. A stationary strategy depends only on
the current state; a positional strategy in a deterministic graph game
chooses one action at each controlled vertex. Nash equilibrium at an initial
state and equilibrium after every history are different requirements.
An expectation of a pathwise $\liminf$ payoff, a $\liminf$ of expected averages,
a limiting expected average, and a uniform finite-horizon equilibrium are
different criteria. None is silently substituted for another.

\textbf{Allocation and mechanisms.} A complete allocation assigns every item
exactly once unless otherwise stated. Zero-valued goods, negative values,
capacity constraints, feasibility, lotteries, and copies of goods affect
fairness definitions and are specified locally. Dominant-strategy incentive
compatibility, Bayesian incentive compatibility, universal truthfulness,
and truthfulness in expectation impose different inequalities. Whenever
an objective ratio has zero denominator, its convention is stated or those
instances are excluded.

\textbf{Infinite-dimensional models.} Expectations, integrals, stochastic
equations, and optimization objectives require their local regularity,
integrability, filtrations, and admissible domains. A backward--forward
mean-field system is a boundary-value problem; it is not an initial-value
semiflow without an additional construction. Long-time, large-population,
and vanishing-noise limits require an order of limits and a convergence
topology. If a minimum need not be attained, the objective is its infimum
or supremum and the approximation requirement is stated separately.
% END ECONOMICS STATEMENT
\end{document}
